\begin{problem}[Indexing with sigma notation]
Given
\[
\sum_{k=1}^{n}(3k-1)=119,
\]
find $n$.
\end{problem}
\begin{hint}
Split the sum into $3\sum k - \sum 1$.
\end{hint}
\begin{solution}
\[
\sum_{k=1}^{n}(3k-1)=3\cdot\frac{n(n+1)}{2}-n=\frac{3n^2+n}{2}.
\]
Set equal to $119$:
\[
\frac{3n^2+n}{2}=119 \implies 3n^2+n-238=0.
\]
\[
\Delta=1+2856=2857.
\]
This is not a perfect square, so there is no integer $n$ with exact value $119$.
\end{solution}

\begin{problem}[Infinite geometric interpretation]
Evaluate
\[
0.4\overline{2}=0.42222\ldots
\]
as a fraction.
\end{problem}
\begin{hint}
Write as $0.4+0.02222\ldots$ and convert the recurring part using a GP.
\end{hint}
\begin{solution}
\[
0.02222\ldots=\frac{2}{90}=\frac{1}{45}.
\]
So
\[
0.4\overline{2}=\frac{2}{5}+\frac{1}{45}=\frac{18}{45}+\frac{1}{45}=\frac{19}{45}.
\]
\end{solution}

\begin{problem}[Second-kind Chebyshev recurrence]
The second-kind Chebyshev polynomials satisfy
\[
U_0(x)=1,\qquad U_1(x)=2x,\qquad U_{n+1}(x)=2xU_n(x)-U_{n-1}(x).
\]
Use this recurrence to find $U_3(x)$ and $U_4(x)$.
\end{problem}
\begin{hint}
First find $U_2(x)$, then continue to $U_3(x)$ and $U_4(x)$.
\end{hint}
\begin{solution}
\[
U_2(x)=2x(2x)-1=4x^2-1.
\]
\[
U_3(x)=2x(4x^2-1)-2x=8x^3-4x.
\]
\[
U_4(x)=2x(8x^3-4x)-(4x^2-1)=16x^4-12x^2+1.
\]
\end{solution}
\begin{takeaways}
\item The first-kind and second-kind Chebyshev polynomials use the same recurrence form.
\item The starting values determine which family is produced.
\item Organize recurrence work line by line to avoid algebra slips.
\end{takeaways}

\begin{problem}[Alternative telescoping patterns]
\begin{enumerate}[(i)]
\item Evaluate $\sum_{n=0}^{4}(-1)^{n+1}(n+5)$.
\item Show $\frac{1}{n(n+1)}=\frac1n-\frac1{n+1}$.
\item Hence find
\[
S_N=\sum_{n=1}^N\left(\frac1n-\frac1{n+1}\right)
\]
and $\lim_{N\to\infty}S_N$.
\end{enumerate}
\end{problem}
\begin{hint}
Part (iii) is a direct telescoping sum after part (ii).
\end{hint}
\begin{solution}
\[
\sum_{n=0}^{4}(-1)^{n+1}(n+5)=-7.
\]
Also
\[
\frac1{n(n+1)}=\frac{(n+1)-n}{n(n+1)}=\frac1n-\frac1{n+1}.
\]
So
\[
S_N=1-\frac1{N+1},\qquad \lim_{N\to\infty}S_N=1.
\]
\end{solution}

\begin{problem}[Surd telescoping]
Let
\[
a_k=\frac1{\sqrt{k+1}+\sqrt{k}},\qquad S_n=\sum_{k=1}^n a_k.
\]
\begin{enumerate}[(i)]
\item Show $a_k=\sqrt{k+1}-\sqrt{k}$.
\item Find $S_n$.
\item Evaluate $S_{99}$.
\item Find the smallest integer $N$ such that $S_N>10$.
\end{enumerate}
\end{problem}
\begin{hint}
Rationalize with the conjugate and then telescope.
\end{hint}
\begin{solution}
\[
a_k=\frac{\sqrt{k+1}-\sqrt{k}}{(k+1)-k}=\sqrt{k+1}-\sqrt{k}.
\]
Hence
\[
S_n=(\sqrt2-\sqrt1)+\cdots+(\sqrt{n+1}-\sqrt n)=\sqrt{n+1}-1.
\]
So $S_{99}=10-1=9$.
For $S_N>10$:
\[
\sqrt{N+1}-1>10 \implies N>120,
\]
smallest such integer is $N=121$.
\end{solution}

\begin{problem}[Central symmetry in a GP]
In a GP, define $P_n=T_1T_2\cdots T_n$.
\begin{enumerate}[(i)]
\item Show $P_n=a^n r^{\frac{n(n-1)}2}$.
\item Show $T_kT_{n-k+1}=T_1T_n$.
\item If $n=11$ and middle term $T_6=5$, find $P_{11}$.
\end{enumerate}
\end{problem}
\begin{hint}
Use $T_k=ar^{k-1}$ and pair symmetric terms.
\end{hint}
\begin{solution}
\[
P_n=a(ar)\cdots(ar^{n-1})=a^n r^{0+1+\cdots+(n-1)}=a^n r^{\frac{n(n-1)}2}.
\]
Also
\[
T_kT_{n-k+1}=ar^{k-1}\cdot ar^{n-k}=a^2r^{n-1}=T_1T_n.
\]
For odd $n=11$, the product is middle-term power:
\[
P_{11}=T_6^{11}=5^{11}.
\]
\end{solution}

\begin{problem}[Fibonacci and Lucas]
Define
\[
F_1=1,\ F_2=1,\ F_n=F_{n-1}+F_{n-2},
\]
\[
L_1=1,\ L_2=3,\ L_n=L_{n-1}+L_{n-2}.
\]
Show that
\[
L_n=F_{n-1}+F_{n+1}.
\]
\end{problem}
\begin{hint}
Set $A_n=L_n+F_n$ and $B_n=L_n-F_n$, then compare with Fibonacci terms.
\end{hint}
\begin{solution}
From initial terms and recurrence,
\[
A_n=L_n+F_n=2F_{n+1},\qquad B_n=L_n-F_n=2F_{n-1}.
\]
Adding gives
\[
2L_n=2F_{n+1}+2F_{n-1},
\]
so
\[
L_n=F_{n-1}+F_{n+1}.
\]
\end{solution}

\begin{problem}[Linear combinations of two GPs]
Let $T_n=ar^{n-1}$ and $U_n=AR^{n-1}$.
Define $W_n=T_n+U_n$.
\begin{enumerate}[(i)]
\item Show $(W_n)$ satisfies
\[
W_{n+2}-(r+R)W_{n+1}+rRW_n=0.
\]
\item Solve
\[
X_{n+2}-5X_{n+1}+6X_n=0,\quad X_1=5,\ X_2=13.
\]
\end{enumerate}
\end{problem}
\begin{hint}
For (ii), the characteristic roots are $2$ and $3$.
\end{hint}
\begin{solution}
Substitute $W_n=ar^{n-1}+AR^{n-1}$ directly to verify (i).
For (ii), write
\[
X_n=\alpha 2^{n-1}+\beta 3^{n-1}.
\]
Using $X_1=5,\ X_2=13$:
\[
\alpha+\beta=5,\quad 2\alpha+3\beta=13
\]
gives $\alpha=2,\ \beta=3$.
Hence
\[
X_n=2^n+3^n.
\]
\end{solution}

\begin{problem}[Partial fractions and telescoping series]

Let 

$$U_n = \frac{1}{n(n+1)(n+2)}$$ 

where $n$ is a positive integer.

\textbf{(i)} Express the general term $U_n$ in partial fractions by decomposing it into the form:
\[
\frac{1}{n(n+1)(n+2)} = \frac{A}{n} + \frac{B}{n+1} + \frac{C}{n+2}
\]
where $A$, $B$, and $C$ are constants to be determined.

\textbf{(ii)} By grouping the resulting fractions to reveal the telescoping nature of the sequence, find an expression for the sum of the first $N$ terms, 

$$S_N = \sum_{n=1}^{N} U_n = \sum_{n=1}^{N} \frac{1}{n(n+1)(n+2)}$$.

\end{problem}

\begin{hint}
\begin{itemize}
    \item \textbf{Part (i):} Multiply both sides of the equation by the common denominator $n(n+1)(n+2)$. To quickly find $A$, $B$, and $C$, substitute values for $n$ that make the bracketed terms equal to zero (e.g., $n=0, -1, -2$).
    \item \textbf{Part (ii) - Grouping:} You will find that $B = -1$. It is helpful to factor out $\frac{1}{2}$ from the entire expression. Then, split the middle $-2$ term into two separate $-1$ fractions: $-\frac{1}{n+1} - \frac{1}{n+1}$. Group one with the $\frac{1}{n}$ term and the other with the $\frac{1}{n+2}$ term.
    \item \textbf{Part (ii) - Summing:} Write out the grouped terms for $n=1$, $n=2$, and $n=N$ explicitly. Look for the diagonal cancellation pattern to see which terms at the beginning and end survive.
\end{itemize}
\end{hint}

\begin{solution}
\textbf{(i) Decompose the general term:} \\
Multiply by the denominator:
\[
1 = A(n+1)(n+2) + B(n)(n+2) + C(n)(n+1)
\]
Substitute values to solve for the constants:
\begin{itemize}
    \item Let $n = 0 \implies 1 = A(1)(2) \implies A = \frac{1}{2}$
    \item Let $n = -1 \implies 1 = B(-1)(1) \implies B = -1$
    \item Let $n = -2 \implies 1 = C(-2)(-1) \implies C = \frac{1}{2}$
\end{itemize}
Therefore:
\[
U_n = \frac{1/2}{n} - \frac{1}{n+1} + \frac{1/2}{n+2} = \frac{1}{2}\left(\frac{1}{n} - \frac{2}{n+1} + \frac{1}{n+2}\right)
\]

\textbf{(ii) Group and sum the sequence:} \\
Group the fractions inside the parenthesis to form differences:
\[
U_n = \frac{1}{2}\left[\left(\frac{1}{n} - \frac{1}{n+1}\right) - \left(\frac{1}{n+1} - \frac{1}{n+2}\right)\right]
\]
Expand $S_N$ to observe the cascading cancellations:
\begin{align*}
    n&=1: \frac{1}{2}\left[\left(\frac{1}{1} - \frac{1}{2}\right) - \left(\frac{1}{2} - \frac{1}{3}\right)\right] \\
    n&=2: \frac{1}{2}\left[\left(\frac{1}{2} - \frac{1}{3}\right) - \left(\frac{1}{3} - \frac{1}{4}\right)\right] \\
    &\dots \\
    n&=N: \frac{1}{2}\left[\left(\frac{1}{N} - \frac{1}{N+1}\right) - \left(\frac{1}{N+1} - \frac{1}{N+2}\right)\right]
\end{align*}
Summing vertically, all intermediate brackets cancel, leaving only the first half of the $n=1$ term and the second half of the $n=N$ term:
\[
S_N = \frac{1}{2}\left[\left(1 - \frac{1}{2}\right) - \left(\frac{1}{N+1} - \frac{1}{N+2}\right)\right] = \frac{1}{2}\left(\frac{1}{2} - \frac{1}{N+1} + \frac{1}{N+2}\right)
\]
\end{solution}

\begin{takeaways}
    \item \textbf{The Method of Cover-up:} When decomposing partial fractions with distinct linear factors, substituting the roots of the denominator is the fastest way to find the numerator constants.
    \item \textbf{Structuring Telescoping Series:} A 3-term partial fraction expansion (like $\frac{1}{2}, -1, \frac{1}{2}$) can often be factored to reveal a $(1, -2, 1)$ pattern, which hides a two-part telescoping series. Splitting the middle coefficient transforms it into a recognizable $V_n - V_{n+1}$ pattern.
    \item \textbf{Write it Out:} Always manually expand the first few and last few terms of a telescoping sum. Relying solely on mental math for the cancellations frequently leads to dropping surviving terms.
\end{takeaways}
